\subsection*{CU-46: Agregar un moderador}
\addcontentsline{toc}{subsection}{CU-46: Agregar un moderador}
\label{cu-46}

\begin{fichacu}{CU-46}{Agregar un moderador}
  \crudFila{Responsable}{Ángel Gabriel Aguilar Hernández}
  \crudFila{Actor principal}{Administrador}
  \crudFila{Actores de sistema}{Servidor de partidas, Base de datos}
  \crudFila{Prototipos}{—}
  \crudFila{Objetivo}{Permitir al Administrador otorgar el rol de moderación a un jugador de confianza, dejando constancia de quién lo otorgó, para que pueda revisar reportes y aplicar sanciones.}
  \crudFila{Disparador}{El Administrador selecciona ``agregar moderador'' en la \texttt{GUI\_\allowbreak{}AdminPanel}.}
  \textbf{Precondiciones} & \cuCelda{\begin{cureglas}
      \item \textbf{\mbox{PRE-01}} El Administrador tiene una \textit{Sesion} abierta y su \textit{Usuario} tiene \texttt{rol} igual a \texttt{ADMINISTRADOR}.
      \item \textbf{\mbox{PRE-02}} El Servidor de partidas está en ejecución y tiene conexión disponible con la Base de datos.
    \end{cureglas}} \\ \hline
  \textbf{Flujo normal} & \cuCelda{\begin{cuenum}[start=1]
      \item El Administrador selecciona ``agregar moderador''.
      \item El SJM despliega un buscador por \texttt{nickname} o \texttt{codigo\_\allowbreak{}amigo}.
      \item El Administrador busca al jugador y lo selecciona de los resultados. \textit{(Ver \mbox{FA-01})}
      \item El SJM despliega la ficha del candidato con su \texttt{nickname}, su \texttt{fecha\_\allowbreak{}registro}, su \texttt{estado\_\allowbreak{}cuenta}, su \texttt{tipo\_\allowbreak{}cuenta}, su historial de \textit{Sancion} y los \textit{Reporte} recibidos, y la opción ``Otorgar rol de moderador''.
      \item El Administrador selecciona la opción, escribe el motivo y confirma en la \texttt{GUI\_\allowbreak{}MessageConfirm}. \textit{(Ver \mbox{FA-02})}
      \item El SJM envía el mensaje \texttt{GRANT\_\allowbreak{}MODERATOR\_\allowbreak{}REQUEST} con el \texttt{id\_\allowbreak{}usuario}, el motivo y el token. \textit{(Ver \mbox{EX-03}, \mbox{EX-04})}
      \item El Servidor de partidas verifica que el token corresponda a una \textit{Sesion} abierta y que el \texttt{rol} del solicitante sea \texttt{ADMINISTRADOR} (\mbox{RN-01}). \textit{(Ver \mbox{FA-03})}
      \item El Servidor de partidas verifica que el candidato tenga \texttt{tipo\_\allowbreak{}cuenta} en \texttt{REGISTRADA} y \texttt{estado\_\allowbreak{}cuenta} en \texttt{ACTIVA} (\mbox{RN-02}). \textit{(Ver \mbox{FA-04})}
    \end{cuenum}} \\
   & \cuCelda{\begin{cuenum}[start=9]
      \item El Servidor de partidas verifica que el candidato no tenga ninguna \textit{Sancion} activa (\mbox{RN-02}). \textit{(Ver \mbox{FA-04})}
      \item El Servidor de partidas verifica que el \texttt{rol} del candidato sea \texttt{JUGADOR}. \textit{(Ver \mbox{FA-05})}
      \item El Servidor de partidas inicia una transacción, cambia el \texttt{rol} del candidato a \texttt{MODERADOR} y registra en la \textit{BitacoraModeracion} la acción \texttt{ROL\_\allowbreak{}OTORGADO} con el administrador, el afectado, el motivo y la fecha. \textit{(Ver \mbox{EX-01}, \mbox{EX-02})}
      \item El Servidor de partidas notifica al candidato, si tiene sesión abierta, con \texttt{ROLE\_\allowbreak{}GRANTED}.
      \item El Servidor de partidas responde con \texttt{GRANT\_\allowbreak{}MODERATOR\_\allowbreak{}OK}.
      \item El SJM muestra la confirmación y añade al jugador a la lista de moderadores.
      \item Fin del caso de uso.
    \end{cuenum}} \\ \hline
  \textbf{Flujos alternos} & \cuCelda{\textbf{\mbox{FA-01}: La búsqueda no encuentra al jugador}\par
    \begin{cuenum}[start=1]
      \item El SJM muestra un estado vacío indicando que no hay jugadores con ese nombre o código.
      \item El flujo continúa en el paso 3 del FN.
    \end{cuenum}} \\
   & \cuCelda{\textbf{\mbox{FA-02}: El Administrador cancela}\par
    \begin{cuenum}[start=1]
      \item El Administrador selecciona ``Cancelar'' en la confirmación.
      \item El SJM cierra la confirmación sin enviar nada.
      \item Fin del caso de uso.
    \end{cuenum}} \\
   & \cuCelda{\textbf{\mbox{FA-03}: La sesión no es válida o el usuario no es administrador}\par
    \begin{cuenum}[start=1]
      \item El Servidor de partidas responde con \texttt{SESSION\_\allowbreak{}INVALID} o con \texttt{FORBIDDEN}.
      \item Si la sesión era válida pero falta el rol, el Servidor de partidas registra el intento en la \textit{BitacoraModeracion}.
      \item El SJM despliega la \texttt{GUI\_\allowbreak{}Login} o el menú principal.
      \item Fin del caso de uso.
    \end{cuenum}} \\
   & \cuCelda{\textbf{\mbox{FA-04}: El candidato no cumple los requisitos}\par
    \begin{cuenum}[start=1]
      \item El Servidor de partidas responde con \texttt{GRANT\_\allowbreak{}MODERATOR\_\allowbreak{}FAILED} y el motivo: \texttt{CUENTA\_\allowbreak{}INVITADO}, \texttt{CUENTA\_\allowbreak{}NO\_\allowbreak{}ACTIVA} o \texttt{SANCION\_\allowbreak{}ACTIVA}.
      \item El SJM muestra la \texttt{GUI\_\allowbreak{}MessageWarning} con el motivo.
      \item Fin del caso de uso.
    \end{cuenum}} \\
   & \cuCelda{\textbf{\mbox{FA-05}: El candidato ya es moderador o administrador}\par
    \begin{cuenum}[start=1]
      \item El Servidor de partidas responde con \texttt{GRANT\_\allowbreak{}MODERATOR\_\allowbreak{}FAILED} y el motivo \texttt{YA\_\allowbreak{}TIENE\_\allowbreak{}ROL}, con el rol actual.
      \item El SJM lo indica en la ficha del candidato.
      \item Fin del caso de uso.
    \end{cuenum}} \\
   & \cuCelda{\textbf{\mbox{FA-06}: El Administrador consulta la lista de moderadores}\par
    \begin{cuenum}[start=1]
      \item El Administrador selecciona ``moderadores'' en la \texttt{GUI\_\allowbreak{}AdminPanel}.
      \item El Servidor de partidas recupera los \textit{Usuario} con \texttt{rol} igual a \texttt{MODERADOR}, con la fecha en que se les otorgó según la \textit{BitacoraModeracion} y el número de reportes que han resuelto.
      \item El SJM despliega la lista, con acceso a retirar el rol en el \mbox{CU-47}.
      \item Fin del caso de uso.
    \end{cuenum}} \\ \hline
  \textbf{Excepciones} & \cuCelda{\textbf{\mbox{EX-01}: Fallo al escribir en la base de datos}\par
    \begin{cuenum}[start=1]
      \item El Servidor de partidas revierte la transacción, de modo que no quede un rol otorgado sin su registro en la bitácora.
      \item El Servidor de partidas registra la excepción con un identificador de incidencia y responde con \texttt{SERVER\_\allowbreak{}ERROR}.
      \item El SJM muestra la \texttt{GUI\_\allowbreak{}MessageError} con el identificador.
      \item Fin del caso de uso.
    \end{cuenum}} \\
   & \cuCelda{\textbf{\mbox{EX-02}: Fallo al confirmar la transacción}\par
    \begin{cuenum}[start=1]
      \item El Servidor de partidas trata la transacción como fallida y registra la excepción señalando que el estado es indeterminado.
      \item El SJM recarga la ficha del candidato, que mostrará su rol real.
      \item Fin del caso de uso.
    \end{cuenum}} \\
   & \cuCelda{\textbf{\mbox{EX-03}: Se agota el tiempo de espera de la respuesta}\par
    \begin{cuenum}[start=1]
      \item El SJM no reenvía la solicitud y recarga la ficha del candidato.
      \item Fin del caso de uso.
    \end{cuenum}} \\
   & \cuCelda{\textbf{\mbox{EX-04}: La conexión se interrumpe}\par
    \begin{cuenum}[start=1]
      \item El SJM muestra la barra de conexión perdida y recarga la ficha al recuperarla.
      \item Fin del caso de uso.
    \end{cuenum}} \\ \hline
  \textbf{Postcondiciones} & \cuCelda{\begin{cureglas}
      \item \textbf{\mbox{POST-01}} Si el caso termina por el FN, el \textit{Usuario} tiene \texttt{rol} igual a \texttt{MODERADOR}.
      \item \textbf{\mbox{POST-02}} Si el caso termina por el FN, la \textit{BitacoraModeracion} registra quién otorgó el rol, a quién, por qué y cuándo.
      \item \textbf{\mbox{POST-03}} Ninguna otra propiedad de la cuenta cambió.
      \item \textbf{\mbox{POST-04}} El nuevo moderador tiene acceso a la cola de moderación en su siguiente solicitud, sin cerrar sesión.
    \end{cureglas}} \\ \hline
  \textbf{Reglas de negocio} & \cuCelda{\begin{cureglas}
      \item \textbf{\mbox{RN-01}} Solo un administrador otorga el rol de moderación.
      \item \textbf{\mbox{RN-02}} Solo puede ser moderador una cuenta \texttt{REGISTRADA}, \texttt{ACTIVA} y sin sanciones activas.
      \item \textbf{\mbox{RN-03}} Los roles posibles son \texttt{JUGADOR}, \texttt{MODERADOR} y \texttt{ADMINISTRADOR}. Un usuario tiene exactamente uno.
      \item \textbf{\mbox{RN-04}} Las cuentas de administración no se crean desde la aplicación: se crean por SQL, igual que los catálogos precargados (\mbox{D-09}).
      \item \textbf{\mbox{RN-05}} Todo cambio de rol queda en la \textit{BitacoraModeracion} con su autor y su motivo.
      \item \textbf{\mbox{RN-06}} El rol se comprueba en el servidor en cada solicitud, así que el cambio surte efecto sin necesidad de volver a iniciar sesión.
    \end{cureglas}} \\ \hline
  \textbf{Observación} & \cuCelda{\textit{Sobre por qué el rol es un atributo y no una tabla.}\par
    Un usuario tiene exactamente un rol y los tres roles son excluyentes, así que \texttt{rol} es un atributo de \textit{Usuario} con un dominio de tres valores. Una tabla de roles con una relación de muchos a muchos solo tendría sentido si un usuario pudiera tener varios roles a la vez o si los permisos de cada rol se configuraran desde la aplicación, y ninguna de las dos cosas está prevista. La historia de los cambios de rol, que sí interesa conservar, está en la \textit{BitacoraModeracion}.} \\ \hline
  \textbf{Datos que toca} & \cuCelda{\begin{cudatos}
      \item \textit{Sesion}, \textit{Usuario}: lee el \texttt{rol} del administrador \textit{(paso 7)}
      \item \textit{Usuario}: lee el candidato por \texttt{nickname} o \texttt{codigo\_\allowbreak{}amigo} \textit{(pasos 3, 8)}
      \item \textit{Sancion}, \textit{Reporte}: lee el historial del candidato \textit{(pasos 4, 9)}
      \item \textit{Usuario}: escribe \texttt{rol} \textit{(paso 11)}
      \item \textit{BitacoraModeracion}: inserta \textit{(pasos 11, \mbox{FA-03})}
    \end{cudatos}} \\ \hline
\end{fichacu}
\clearpage
